\documentclass{article}
\usepackage[T1]{fontenc}
\usepackage{amsmath,amssymb,amsthm}
\usepackage[margin=1in]{geometry}

\newtheorem{proposition}{Proposition}

\begin{document}

\begin{center}
  \Large\bfseries The leading logarithmic law for the effective coupling
\end{center}

Let \(G\subset U(r_G)\) be a closed matrix subgroup, in the indicated
unitary representation, with nonabelian simple real Lie algebra
\(\mathfrak g\). Write
\[
    b_0(G)=\frac{11\,C_{\mathfrak g}}{48\pi^2},
    \qquad
    C_{\mathfrak g}
      =-\frac{1}{\dim_{\mathbb R}\mathfrak g}
        \sum_j \operatorname{Tr}_{\mathfrak g}
             \bigl((\operatorname{ad}e_j)^2\bigr),
\]
where \((e_j)\) is an orthonormal basis of \(\mathfrak g\) for
\(\langle X,Y\rangle=-r_G^{-1}\operatorname{Re}\operatorname{tr}(XY)\).
Here \(\operatorname{Tr}_{\mathfrak g}\) is the real trace of an
endomorphism of \(\mathfrak g\). Invariance of the matrix trace gives
\(\langle[X,Y],Z\rangle=-\langle Y,[X,Z]\rangle\), so each
\(\operatorname{ad}e_j\) is skew-adjoint. Consequently
\(-\operatorname{Tr}_{\mathfrak g}((\operatorname{ad}e_j)^2)
=\sum_\ell\|[e_j,e_\ell]\|^2\). Noncommutativity of \(\mathfrak g\)
makes at least one term in the resulting double sum positive. In particular,
\(b_0(G)>0\).

\paragraph{Input from the lattice construction.}
Fix a construction family \(\mathcal C\) satisfying the construction
conditions. For \(N,p,m\in\mathbb N\) with \(m\le N\), let
\(\beta_{N,p,m}\) be its action coefficient at stage \(m\) of the
trajectory started at lattice level \(N\) and period index \(p\).
The spacing at stage \(m\) is \(a_m=(\prod_{j<m}B_j)^{-1}\), where each
block factor \(B_j\) is an integer at least \(3\). Thus
\(s_m=\log(a_m^{-1})=\sum_{j<m}\log B_j\ge m\log 3\), so \(s_m\to+\infty\).
The running estimate for the family provides a fixed reference intercept
\(\bar\beta\) and, for every
\(\eta>0\), a threshold \(m_0(\eta)\) such that
\begin{equation}
 \left|\beta_{N,p,m}-\bigl(\bar\beta+2b_0(G)s_m\bigr)\right|
 \leq \eta s_m
 \qquad
 \bigl(m\geq m_0(\eta),\ N\geq m,\ p\in\mathbb N\bigr).
 \label{eq:running}
\end{equation}
The threshold is uniform over the starting level and period index.
We use the running estimate as an input to the comparison below.

\begin{proposition}
Let \(N,k:\mathbb N\to\mathbb N\) be arbitrary sequences with \(i\leq N(i)\)
for all sufficiently large \(i\). For \(\beta_{N,p,m}>0\) set
\(g_{N,p,m}=(\sqrt{\beta_{N,p,m}})^{-1}\). Then
\(\beta_{N(i),k(i),i}>0\) for all sufficiently large \(i\), and
\[
  \lim_{i\to\infty}g_{N(i),k(i),i}\sqrt{\log(a_i^{-1})}
  =\frac{1}{\sqrt{2b_0(G)}}.
\]
\end{proposition}

\begin{proof}
Fix \(\varepsilon>0\). Use \eqref{eq:running} with
\(\eta=\varepsilon/2\), and take \(m\) large enough also to have \(s_m>0\)
and \(\lvert\bar\beta\rvert/s_m<\varepsilon/2\). For every \(N\ge m\)
and every \(p\), the triangle inequality yields
\begin{align*}
 \left|\frac{\beta_{N,p,m}}{s_m}-2b_0(G)\right|
 &\leq \frac{\left|\beta_{N,p,m}
              -\bigl(\bar\beta+2b_0(G)s_m\bigr)\right|}{s_m}
       +\frac{|\bar\beta|}{s_m} \\
 &< \varepsilon .
\end{align*}
The same threshold therefore works at \(m=i\), \(N=N(i)\), \(p=k(i)\)
for all sufficiently large \(i\). It follows that
\(\beta_{N(i),k(i),i}/s_i\to 2b_0(G)>0\), which also gives the asserted
eventual positivity. On that tail both factors under square roots are
positive and
\[
 g_{N(i),k(i),i}\sqrt{\log(a_i^{-1})}
 =\left(\frac{\beta_{N(i),k(i),i}}{s_i}\right)^{-1/2}.
\]
The map \(x\mapsto x^{-1/2}\) is continuous on \((0,\infty)\), giving
the limit.
\end{proof}

At \(N(i)=i\) the coefficient is the bare lattice action coefficient.
For a general admissible \(N(i)\), the comparison concerns the coefficient
read from the trajectory at spacing \(a_i\). Relating this inverse-square-root
proxy to a coupling defined through a physical observable would require a
separate matching theorem.

\end{document}
